% ============================================================
%  PROBLEMAS DE PROBABILIDAD — UNIDAD 1
%  Un archivo por unidad; agregar bloques \dayheader según avanza el curso.
%  Compilar con pdflatex (funciona en Overleaf sin configuración extra).
% ============================================================

\documentclass[12pt, a4paper]{article}

% ---- codificación e idioma ---------------------------------
\usepackage[utf8]{inputenc}
\usepackage[T1]{fontenc}
\usepackage[spanish]{babel}

% ---- matemática --------------------------------------------
\usepackage{amsmath, amssymb, amsthm}

% ---- código R ----------------------------------------------
\usepackage{listings}
\usepackage{xcolor}

\definecolor{codegreen}{rgb}{0.2,0.5,0.2}
\definecolor{codegray}{rgb}{0.5,0.5,0.5}
\definecolor{codeblue}{rgb}{0.1,0.1,0.8}
\definecolor{backcolor}{rgb}{0.97,0.97,0.97}

\lstdefinestyle{Rstyle}{
  language=R,
  backgroundcolor=\color{backcolor},
  basicstyle=\ttfamily\small,
  keywordstyle=\color{codeblue}\bfseries,
  commentstyle=\color{codegreen}\itshape,
  stringstyle=\color{orange},
  numberstyle=\tiny\color{codegray},
  numbers=left,
  stepnumber=1,
  frame=single,
  rulecolor=\color{codegray},
  breaklines=true,
  showstringspaces=false,
  tabsize=2,
  captionpos=b
}
\lstset{style=Rstyle}

% ---- diseño de página --------------------------------------
\usepackage{geometry}
\geometry{margin=2.5cm}

\usepackage{enumitem}
\usepackage{parskip}
\usepackage{titlesec}
\usepackage{fancyhdr}
\usepackage{lastpage}
\usepackage{hyperref}

% ---- encabezado y pie de página ----------------------------
\pagestyle{fancy}
\fancyhf{}
\lhead{\textbf{Probabilidad} — Unidad 1: \nombreunidad}
\rhead{\today}
\cfoot{Página \thepage\ de \pageref{LastPage}}
\renewcommand{\headrulewidth}{0.4pt}

% ---- contador de problemas ---------------------------------
%  Se reinicia con cada \dayheader.
\newcounter{problema}
\renewcommand{\theproblema}{\arabic{diaactual}.\arabic{problema}}

\newcounter{diaactual}   % guarda el número del día actual

\newcommand{\problema}[1][]{%        \problema  o  \problema[título]
  \refstepcounter{problema}%
  \medskip
  \noindent\textbf{Problema~\theproblema%
    \ifx&#1&\else~—~#1\fi}%
  \par\smallskip
}

% ---- comando \dayheader ------------------------------------
%  Uso:  \dayheader{1}{Subtítulo descriptivo}
%  (Se llama dayheader en lugar de \day para no colisionar
%   con el comando interno de LaTeX \day.)
\newcommand{\dayheader}[2]{%
  \setcounter{diaactual}{#1}%
  \setcounter{problema}{0}%
  \section*{Clase #1 \normalfont\small— #2}%
  \addcontentsline{toc}{section}{Clase #1: #2}%
}

% ---- nombre de la unidad (editar aquí) ---------------------
\newcommand{\nombreunidad}{Espacios Muestrales y Eventos}

% ============================================================
\begin{document}
% ============================================================

\begin{center}
  {\LARGE \textbf{Unidad 1: \nombreunidad}} \\[0.4em]
  {\large Probabilidad — Guía de Problemas} \\[0.2em]
  {\small Última actualización: \today}
\end{center}

\tableofcontents
\bigskip\hrule\bigskip

% ============================================================
\dayheader{1}{Introducción a los espacios muestrales}
% ============================================================

\problema
Se lanza un dado equilibrado de seis caras.
\begin{enumerate}[label=(\alph*)]
  \item Escriba el espacio muestral $\Omega$.
  \item Sea $A$ el evento ``el resultado es par''. Liste los elementos de $A$.
  \item Calcule $P(A)$.
\end{enumerate}

\problema[Regla del complemento]
Sea $B$ el evento ``el resultado es mayor que 4''.
\begin{enumerate}[label=(\alph*)]
  \item Calcule $P(B)$ y $P(B^c)$.
  \item Verifique que $P(B) + P(B^c) = 1$.
\end{enumerate}

\problema[Simulación en R]
El siguiente código simula 10\,000 lanzamientos de un dado equilibrado
y estima $P(A)$ del Problema~1.1.
\begin{lstlisting}
set.seed(42)
lanzamientos <- sample(1:6, size = 10000, replace = TRUE)

# Estimar P(par)
mean(lanzamientos %% 2 == 0)
\end{lstlisting}
Ejecute el código, compare con el valor teórico y explique cualquier discrepancia.

% ============================================================
\dayheader{2}{Operaciones entre eventos}
% ============================================================

\problema
Sea $\Omega = \{1,2,3,4,5,6\}$, $A = \{2,4,6\}$, $B = \{4,5,6\}$.
Calcule:
\begin{enumerate}[label=(\alph*)]
  \item $A \cup B$
  \item $A \cap B$
  \item $A \setminus B$
  \item $A^c \cap B$
\end{enumerate}

\problema
Con los mismos eventos, verifique la fórmula de inclusión-exclusión:
\[
  P(A \cup B) = P(A) + P(B) - P(A \cap B).
\]

\problema[Leyes de De Morgan]
Demuestre o ilustre con un ejemplo que
\[
  (A \cup B)^c = A^c \cap B^c.
\]

% ============================================================
\dayheader{3}{Probabilidad condicional}
% ============================================================

\problema
Se extrae una carta al azar de un mazo estándar de 52 cartas.
Sea $A = \{\text{la carta es corazón}\}$ y $B = \{\text{la carta es figura}\}$.
\begin{enumerate}[label=(\alph*)]
  \item Calcule $P(A \mid B)$.
  \item ¿Son $A$ y $B$ independientes? Justifique.
\end{enumerate}

\problema[Simulación en R]
Estime $P(A \mid B)$ por simulación y compare con el resultado anterior.
\begin{lstlisting}
set.seed(7)
palos   <- rep(c("corazon","diamante","trebol","pica"), each = 13)
valores <- rep(c(2:10, "J", "Q", "K", "A"), times = 4)
mazo    <- data.frame(palo = palos, valor = valores)

n       <- 100000
extracc <- mazo[sample(nrow(mazo), n, replace = TRUE), ]

figuras  <- extracc$valor %in% c("J", "Q", "K")
corazon  <- extracc$palo == "corazon"

# Estimar P(corazon | figura)
mean(corazon[figuras])
\end{lstlisting}

% ============================================================
%  Agregar más bloques \dayheader{4}{...} a medida que avanza la unidad.
% ============================================================

\end{document}